\documentclass[12pt]{article}

% \usepackage{a4wide}
\thispagestyle{empty}
\begin{document}
%\newfont{\twlsy}{cmsy10 scaled 1200}

\Large
{
\begin{center}
Analysing the Thermodynamics of Steam Reforming with Maple
\end{center}
}

{\small
\begin{center} Johannes Grotendorst${}^*$ and J\"urgen Dornseiffer${}^\dag$ \\
\medskip
${}^*$ Central Institute for Applied Mathematics\\
${}^\dag$ Institute for Applied Physical Chemistry\\
Research Centre J\"ulich\\  D-52425 J\"ulich, Germany\\
\{j.grotendorst, j.dornseiffer\}@fz-juelich.de
\end{center}
}

\begin{abstract}
The catalytic steam reforming process converts hydrocarbons into mixtures of
hydrogen, carbon monoxide, carbon dioxide, and methane. It is widely
used in industry for the production of hydrogen, synthesis gas, ammonia,
methanol and other useful chemicals.
Important data for the design and optimization
of a steam reformer are provided by thermodynamic calculations.
For instance, to find the most economic reaction conditions and for
predicting the carbon boundaries, it is necessary to study
the reaction behavior with respect to the
operating parameters. Steam reforming of a general hydrocarbon
is based on the following simple system of reaction equations:

\medskip

\noindent
1. The methane-steam equality
%\chemical{CH_{4}+H_{2}O \rightleftharpoons CO+3H_{2}} \hfill (1)
$${\rm CH_{4}+H_{2}O} \rightleftharpoons {\rm CO+3H_{2}}$$
2. The water-gas shift
$${\rm CO+H_{2}O} \rightleftharpoons {\rm CO_{2}+H_{2}}$$
3. The carbon-methanation equality
$${\rm CH_{4}} \rightleftharpoons {\rm C+2H_{2}}$$

\medskip
\noindent

Using symbolic methods we introduce a flexible computation model
for the thermodynamic description of the steam reforming process.
The product gas composition at equilibrium is determined by solving
a system of multiparameter nonlinear equations for the reaction extents.
The applied computation techniques are chosen graduated according to
the computation effort needed and range from direct calculations
in Maple to the automatic generation of Fortran programs for complete
parameter studies.
We analyse the effects of temperature, pressure and the
steam to carbon ratio on the equilibrium gas composition and
the carbon formation.
The theoretical equilibrium compositions are found to be closely
matching with those obtained in experiments.
The computed data were used to determine the mass and energy balances
of a pilot reactor for reforming organic waste (fluids and gases).

This work presents an example of interdisciplinary work done
in the field of  scientific computing
at the institutes for applied mathematics and physical chemistry
of the research centre J\"ulich . It illustrates how the use of modern
computing tools such as Maple changes the approach of mathematical
modeling in science. The concept of scientific computation has evolved
from a field encompassing primarily numerical methods to a much broader
field that includes algebraic and analytical methods, numerical methods,
and graphics.

\end{abstract}
\end{document}
